% !PRINTED 46
\weberpage{46}

Unter diesen Charakteren kommt auch der vor, den man erhält, wenn man $\beta_1 = 0, \beta_2 = 0 \ldots \beta_\nu = 0$ setzt, der für alle Elemente $A$ den Werth 1 hat. Diesen wollen wir den \gesperrt{Einheitscharakter} nennen.

Die $n$ Charaktere von $S$ können nun selbst wieder zu einer \gesperrt{Abel}'schen Gruppe vereinigt werden, und zwar zu einer mit $S$ isomorphen Gruppe.

Nach (13) ist nämlich jeder der $n$ Charaktere durch ein nach dem Modulsysteme $a_1, a_2 \ldots a_\nu$ genommenes Zahlensystem $\beta_1, \beta_2 \ldots \beta_\nu$ bestimmt, und dies Zahlensystem der $\beta$ ist zugleich das System der Indices eines bestimmten Elementes $B$ von $S$, nämlich von
\weq{14}{B = A_1^{\beta_1}A_2^{\beta_2}\ldots A_\nu^{\beta_\nu},}
so dass jedem Elemente $B$ von $S$ ein bestimmter Charakter entspricht, den wir mit $\chi_B(A)$ oder, indem wir $A$ weglassen, mit $\chi_B$ bezeichnen. Diese Zuordnung der Charaktere $\chi$ zu den Elementen $B$ wird sich aber ändern, wenn eine andere Basis zu Grunde gelegt wird, oder wenn die Einheitswurzeln $\varepsilon_1, \varepsilon_2 \ldots \varepsilon_\nu$ anders angenommen werden.

Ist $B'$ ein zweites Element von $S$ mit den Indices $\beta_1', \beta_2' \ldots \beta_\nu'$, so erhalten wir in gleicher Weise den Charakter $\chi_{B'}$, und wenn wir nun unter $\chi_{BB'}$ den Charakter verstehen, der für jedes $A$ durch die Formel
\weq{15}{\chi_{BB'}(A) = \varepsilon_1^{\alpha_1(\beta_1+\beta_1')}\varepsilon_2^{\alpha_2(\beta_2+\beta_2')}\ldots\varepsilon_\nu^{\alpha_\nu(\beta_\nu+\beta_\nu')}}
bestimmt ist, so sind die Charaktere $\chi$ hierdurch zu einer mit $S$ isomorphen Gruppe verbunden. Das Einheitselement in dieser \gesperrt{Charakterengruppe} ist der Einheitscharakter.

Es lässt sich hiernach auch die Gruppe der Charaktere durch eine Basis darstellen, die man erhält, wenn man
\weq{16}{\chi_1(A) = \varepsilon_1^{\alpha_1},\ \chi_2(A) = \varepsilon_2^{\alpha_2} \ldots \chi_\nu(A) = \varepsilon_\nu^{\alpha_\nu}}
setzt. Dann ist jeder Charakter in der Form enthalten:
\weq{17}{\chi_B = \chi_1^{\beta_1}\chi_2^{\beta_2}\ldots\chi_\nu^{\beta_\nu},}
und es ist, wenn $B, B'$ zwei Elemente in $S$ sind:
\weq{18}{\chi_B\chi_{B'} = \chi_{BB'}.}

Sind $\chi_1, \chi_2$ irgend zwei der Charaktere und $\chi_1\chi_2$ der aus beiden zusammengesetzte, so ist nach (15) für jedes Element $A$
\weq{19}{\chi_1(A)\chi_2(A) = \chi_1\chi_2(A).}
